% =====================================================================
%  2. Literature Review
% =====================================================================
\section{Literature Review}
\label{sec:literature}

% TODO (before submission): re-check every number below against the original tables
% (several were read through an extraction tool). Sources marked "abstract" in the comments were
% only read at abstract level: read at least Charlton & Singleton, Lloyd and Bouktif in full.

System-level load forecasting methods fall into five groups: classical time-series models,
regression with hand-crafted features, tree ensembles, pattern-based methods, and deep learning.
Results across papers often look contradictory. Most of these differences come from three
factors: the forecast horizon, whether past load is an input, and how weather is provided.

\subsection{Classical statistical baselines}
\label{sec:lit-classical}

Seasonal naive copies the load of the same hour one day or one week earlier. ARIMA (autoregressive integrated moving
average) and exponential smoothing (ETS) model the series itself. On 35 ENTSO-E countries in 2018, Smyl et al.\ report a day-ahead
mean absolute percentage error (MAPE) of 5.08\% for weekly naive, 3.30\% for ARIMA and 3.11\% for ETS~\cite{smyl2023esdrnn}. In
Portugal, weekly seasonal naive reaches 4.29\%~\cite{pelekis2023comparative}. Wang et al.\ conclude
that naive models are not useful as the main benchmark at the aggregate level~\cite{wang2016recency}.
We therefore keep seasonal naive only as a floor (B0).

\subsection{Regression with feature engineering}
\label{sec:lit-regression}

Linear and additive models with calendar, lagged load and temperature terms are the classic
approach. Fan and Hyndman use a semi-parametric additive model with calendar, lagged load and
temperature for the Australian market~\cite{fan2012semiparametric}. The winner of the Global Energy
Forecasting Competition 2012 (GEFCom2012) was a
refined multiple regression with separate models per zone, hour, season and day
type~\cite{charlton2014refined}. In Italy, a regularised linear model using only the previous day's
load is competitive with the TSO forecast once special days are excluded~\cite{incremona2021italian}.
For Spain, Pardo et al.\ find that heating and cooling degree days (HDD, CDD) and seasonality stay significant
after autoregressive terms, on daily data~\cite{pardo2002temperature}. These models are easy to
explain, but they depend on the analyst choosing the right functional form.

\subsection{Tree ensembles and gradient boosting}
\label{sec:lit-trees}

Tree ensembles learn thresholds and interactions themselves. In GEFCom2012, an ensemble that
included gradient boosting machines with little domain tuning ranked second~\cite{lloyd2014gefcom}.
Rubattu et al.\ use LightGBM with lagged and rolling temperature statistics as the most important
inputs; they note that gradient boosting cannot extrapolate a long-term trend, a problem in their
year-ahead round~\cite{rubattu2023electricity}. In Dang-Ha et al., a random forest with lagged load
ranked second on both datasets and was more stable across them than the best
method~\cite{dangha2017local}.

\paragraph{Regression versus boosting.}
The literature we read does not settle which is better. A feature-engineered regression won
GEFCom2012 while a boosting ensemble came second~\cite{hong2013gefcomwinners}, but a competition rank
is not a significance test, and the setting (US zones, backcasting and one-week forecasts) differs
from ours. We found no readable study comparing the two on hourly national load, day-ahead, with
weather. This is exactly the comparison of Q1 and Q2.

\subsection{Pattern-based methods and deep learning}
\label{sec:lit-dl}

Deep learning results are mixed. A meta-regression over 59 studies finds long short-term memory (LSTM) and hybrid networks
best overall~\cite{hopf2023meta}, and an LSTM beat tuned machine-learning models on French
load~\cite{bouktif2018optimal}. On the other hand, LSTM (2.76\% MAPE) was beaten by simple
pattern-based methods on 35 countries~\cite{smyl2023esdrnn}, and in Portugal a plain multilayer perceptron (MLP, 2.02\%)
was close to N-BEATS (Neural Basis Expansion Analysis for Time Series, 1.90\%) at much lower cost~\cite{pelekis2023comparative}. Pattern-based
methods~\cite{dudek2021ensembles} are cheap and strong, but they use no weather, so they cannot
answer Q3. We do not implement LSTM. This is a scope decision: the assignment allows classical
machine learning only, and our questions are about features and weather. It is not a claim that
deep learning is worse.

\subsection{Temperature, information sets and explanation}
\label{sec:lit-weather}

Load is U-shaped in temperature (heating when cold, cooling when hot). Recent temperature history
(lags and moving averages) improved forecasts by 18--21\% over a benchmark in
GEFCom2012~\cite{wang2016recency}. Temperature is strongly collinear with season and hour of
day~\cite{pelekis2023comparative}, so its own effect is hard to isolate.

Hyndman and Athanasopoulos define \emph{ex-ante} forecasts as those using only information
available in advance, and \emph{ex-post} forecasts as those using later information on the
predictors~\cite{hyndman2021fpp3}. Many academic studies use ex-post weather. This is why we report
the \emph{actual} and \emph{lagged} settings separately (\cref{sec:weather-settings}).

For explanation we use SHAP (SHapley Additive exPlanations)
values~\cite{lundberg2017unified}. TreeSHAP, the SHAP algorithm for tree ensembles, only partly accounts for
dependence between features, and explanations can be misleading when features are strongly
dependent~\cite{aas2021explaining}. SHAP explains what the \emph{model} uses, not how load depends
on its causes~\cite{chen2020true}. We therefore read SHAP by feature group, not by single feature.

\subsection{Summary and implications}
\label{sec:lit-summary}

\begin{table}[H]
  \centering
  \small
  \caption{Main lessons from the literature and how they shape our design.}
  \label{tab:lit-summary}
  \begin{tabularx}{\linewidth}{@{}Y Y@{}}
    \toprule
    Finding & Consequence for this work \\
    \midrule
    Seasonal naive is weak but standard~\cite{wang2016recency,smyl2023esdrnn}
      & B0 is a floor; the main opponent of LightGBM is a feature-engineered Ridge. \\
    Regression vs.\ boosting is unresolved~\cite{charlton2014refined,lloyd2014gefcom}
      & No expected winner is assumed; both outcomes are reported with confidence intervals. \\
    Degree days work for Spain~\cite{pardo2002temperature}; recency matters~\cite{wang2016recency}
      & Ridge gets HDD/CDD of current and 24-hour mean temperature. \\
    Boosting cannot extrapolate trends~\cite{rubattu2023electricity}
      & We check bias on validation and allow a ``change since last week'' target. \\
    Most studies use ex-post weather~\cite{hyndman2021fpp3}
      & Two weather settings, reported as optimistic and pessimistic bounds. \\
    Whether a TSO forecast can be beaten depends on the metric and on special-day
      handling~\cite{mobius2023better,incremona2021italian}
      & The TSO column is only a reference; we do not claim to beat or lose to it. \\
    SHAP is unreliable for dependent features~\cite{aas2021explaining,chen2020true}
      & Grouped SHAP, read as ``what the model uses''. \\
    \bottomrule
  \end{tabularx}
\end{table}
